\begin{afterword}
\summaryhead

The post answers a commenter, Cihan Baran, who said that since no situation could make 2 + 2 = 4 false, the belief may be held unconditionally. The author admits being unable to conceive of such a situation too, but says this does not make the belief unconditional.

The author then describes a morning on which two earplugs set beside two others make three, a mental picture of the sum agrees, and calculators, Google and even \textsc{1984} agree too. Faced with that, the author would conclude that 2 + 2 = 3 and explain the old belief as a memory fault or tampering. Such evidence, the post says, is the same kind that now supports 2 + 2 = 4: physical observation, mental visualization and social agreement. A belief that did not get there through a process entangled with reality could be right only by coincidence. So unconditional facts do not need unconditional beliefs. A footnote asks Christians what would convince them of Islam.

\responsehead

The post's first move is correct, and it answers the commenter. That no situation could make a fact false does not mean that nothing could change one's belief in it.

The post then adds a second, larger claim, and states it as plain fact: the evidence for 2 + 2 = 4 is ``exactly the same kind'' as evidence about earplugs. That is a position in the philosophy of mathematics, close to John Stuart Mill's, and in 1884 Frege derided accounts of this kind as ``gingerbread and pebble'' arithmetic. Its arguments for it are the ones the standard objections address. How the author first learned arithmetic does not show what justifies it; Frege placed the a priori in justification, not in how a belief was acquired. Earplugs that add to three would show something about earplugs; the Stanford Encyclopedia's example of a failed sum, liquids that do not add to four quarts, counts against the claim about liquids. The mental picture of xx and xx is close to Kant's counting of units, which Kant used to explain how arithmetic is known a priori. Kant kept the necessity of arithmetic and denied that it comes from experience; Mill kept experience and gave up the necessity. The post keeps both, a fact true in every situation and known from experience, without saying how experience could show the first. Even Quine, who held that arithmetic could in principle be revised, held that no particular experience confirms 2 + 2 = 4.

The larger claim is also not needed. On the a priori view, belief in a proof can be undercut by evidence that one is mad, and the explanations the post reaches for, a memory fault or tampering, are evidence of that kind. The reply to the commenter stands without the empiricism. Three and a half months later, in ``Infinite Certainty'', the author still called 2 + 2 = 4 ``mathematical as well as empirical''.

The commenter's other objection, that a mathematical theorem ``doesn't cause anything'' and so cannot fit a causal definition of evidence, is left unquoted. It is a form of Benacerraf's problem, and the post answers it only by treating arithmetic as a fact about objects.

The footnote to Christians changes the question. For the author, ``what would convince me'' meant evidence; for the Christian reader, the post supposes, it means being raised Muslim.

In short: a correct reply to the commenter, that a necessary truth need not be believed unconditionally, joined to an empiricism about arithmetic that the reply does not need, that the post states as plain fact, and whose arguments Frege, Kant and even Quine rejected.
\end{afterword}
